% ============================================================================
% CAPÍTULO 14: DESIGN OF BIOLOGICAL SYSTEMS
% Bioinformática para Biologia de Sistemas
% ============================================================================

% ============================================================================
% TEMPLATE BASE PARA APRESENTAÇÕES EM BEAMER
% Bioinformática para Biologia de Sistemas
% ============================================================================

\documentclass[12pt, aspectratio=169]{beamer}

% ============================================================================
% PACOTES ESSENCIAIS
% ============================================================================
\usepackage[utf8]{inputenc}
\usepackage[T1]{fontenc}
\usepackage[brazilian]{babel}
\usepackage{amsmath, amssymb, amsthm}
\usepackage{graphicx}
\usepackage{xcolor}
\usepackage{tikz}
\usetikzlibrary{arrows.meta, shapes, positioning, calc, decorations.pathreplacing, decorations.shapes, decorations.pathmorphing}
\usepackage{listings}
\usepackage{booktabs}
\usepackage{multirow}
\usepackage{hyperref}
\usepackage{chemfig}
\usepackage{chemarr}

% ============================================================================
% CONFIGURAÇÃO DO TEMA
% ============================================================================
\usetheme{Madrid}
\usecolortheme{default}

% Cores personalizadas
\definecolor{azulescuro}{RGB}{0, 51, 102}
\definecolor{azulclaro}{RGB}{51, 153, 255}
\definecolor{verdeacido}{RGB}{102, 204, 0}
\definecolor{laranjacel}{RGB}{255, 153, 51}
\definecolor{roxodna}{RGB}{153, 51, 255}

\setbeamercolor{structure}{fg=azulescuro}
\setbeamercolor{palette primary}{bg=azulescuro,fg=white}
\setbeamercolor{palette secondary}{bg=azulclaro,fg=white}
\setbeamercolor{palette tertiary}{bg=azulescuro,fg=white}
\setbeamercolor{palette quaternary}{bg=azulclaro,fg=white}
\setbeamercolor{block title}{bg=azulescuro,fg=white}
\setbeamercolor{block body}{bg=azulclaro!10,fg=black}
\setbeamercolor{block title alerted}{bg=red!80,fg=white}
\setbeamercolor{block body alerted}{bg=red!10,fg=black}
\setbeamercolor{block title example}{bg=verdeacido!80,fg=white}
\setbeamercolor{block body example}{bg=verdeacido!10,fg=black}

% ============================================================================
% CONFIGURAÇÃO DE CÓDIGO
% ============================================================================
\lstset{
    language=Python,
    basicstyle=\ttfamily\small,
    keywordstyle=\color{azulescuro}\bfseries,
    commentstyle=\color{verdeacido},
    stringstyle=\color{laranjacel},
    numbers=left,
    numberstyle=\tiny\color{gray},
    stepnumber=1,
    numbersep=5pt,
    backgroundcolor=\color{white},
    showspaces=false,
    showstringspaces=false,
    showtabs=false,
    frame=single,
    rulecolor=\color{black},
    tabsize=2,
    captionpos=b,
    breaklines=true,
    breakatwhitespace=false,
    escapeinside={\%*}{*)},
    xleftmargin=10pt,
    xrightmargin=10pt
}

% Definir estilo para R
\lstdefinestyle{Rstyle}{
    language=R,
    basicstyle=\ttfamily\small,
    keywordstyle=\color{azulescuro}\bfseries,
    commentstyle=\color{verdeacido},
    stringstyle=\color{laranjacel}
}

% Definir estilo para MATLAB
\lstdefinestyle{MATLABstyle}{
    language=Matlab,
    basicstyle=\ttfamily\small,
    keywordstyle=\color{azulescuro}\bfseries,
    commentstyle=\color{verdeacido},
    stringstyle=\color{laranjacel}
}

% ============================================================================
% COMANDOS PERSONALIZADOS
% ============================================================================

% Comando para equações destacadas
\newcommand{\destaque}[1]{%
    \begin{alertblock}{}
        \begin{center}
            #1
        \end{center}
    \end{alertblock}
}

% Comando para definições
\newcommand{\definicao}[2]{%
    \begin{block}{Definição: #1}
        #2
    \end{block}
}

% Comando para exemplos
\newcommand{\exemplo}[2]{%
    \begin{exampleblock}{Exemplo: #1}
        #2
    \end{exampleblock}
}

% Comando para observações importantes
\newcommand{\importante}[1]{%
    \begin{alertblock}{Importante!}
        #1
    \end{alertblock}
}

% Comandos matemáticos úteis
\newcommand{\R}{\mathbb{R}}
\newcommand{\N}{\mathbb{N}}
\newcommand{\Z}{\mathbb{Z}}
\newcommand{\dif}{\mathrm{d}}
\newcommand{\parcial}[2]{\frac{\partial #1}{\partial #2}}
\newcommand{\derivada}[2]{\frac{\dif #1}{\dif #2}}
\newcommand{\vect}[1]{\boldsymbol{#1}}

% Comandos biológicos
\newcommand{\gene}[1]{\textit{#1}}
\newcommand{\proteina}[1]{\textsc{#1}}
\newcommand{\especie}[1]{\textit{#1}}

% ============================================================================
% CONFIGURAÇÃO DE RODAPÉ
% ============================================================================
\setbeamertemplate{footline}{
  \leavevmode%
  \hbox{%
  \begin{beamercolorbox}[wd=.33\paperwidth,ht=2.25ex,dp=1ex,center]{author in head/foot}%
    \usebeamerfont{author in head/foot}\insertshortauthor
  \end{beamercolorbox}%
  \begin{beamercolorbox}[wd=.34\paperwidth,ht=2.25ex,dp=1ex,center]{title in head/foot}%
    \usebeamerfont{title in head/foot}\insertshorttitle
  \end{beamercolorbox}%
  \begin{beamercolorbox}[wd=.33\paperwidth,ht=2.25ex,dp=1ex,right]{date in head/foot}%
    \usebeamerfont{date in head/foot}\insertshortdate{}\hspace*{2em}
    \insertframenumber{} / \inserttotalframenumber\hspace*{2ex}
  \end{beamercolorbox}}%
  \vskip0pt%
}

% ============================================================================
% CONFIGURAÇÃO DE NAVEGAÇÃO
% ============================================================================
\setbeamertemplate{navigation symbols}{}

% ============================================================================
% CONFIGURAÇÃO DE ITEMIZE/ENUMERATE
% ============================================================================
\setbeamertemplate{itemize items}[circle]
\setbeamertemplate{enumerate items}[default]

% ============================================================================
% FIM DO TEMPLATE
% ============================================================================


% ============================================================================
% INFORMAÇÕES DO DOCUMENTO
% ============================================================================
\title[Cap. 14: Design of Biological Systems]{Capítulo 14: Design of Biological Systems}
\subtitle{Bioinformática para Biologia de Sistemas}
\author{Ney Lemke}
\institute{DBF-Unesp}
\date{\today}

\begin{document}

% ============================================================================
% SLIDE DE TÍTULO
% ============================================================================
\begin{frame}
\titlepage
\end{frame}

% ============================================================================
% SUMÁRIO
% ============================================================================
\begin{frame}{Sumário}
\tableofcontents
\end{frame}

% ============================================================================
% SEÇÃO 1: INTRODUÇÃO
% ============================================================================
\section{Introdução}

\begin{frame}{Visão Geral do Capítulo}
\begin{block}{Objetivos de Aprendizagem}
\begin{itemize}
    \item Compreender princípios de engenharia aplicados à biologia
    \item Projetar e analisar circuitos genéticos sintéticos
    \item Aplicar engenharia metabólica para produção de compostos
    \item Utilizar ferramentas computacionais de design
    \item Entender desafios e perspectivas da biologia sintética
\end{itemize}
\end{block}

\vspace{0.3cm}

\begin{exampleblock}{Biologia Sintética}
Aplicação de princípios de engenharia para projetar e construir novos sistemas biológicos ou redesenhar sistemas naturais existentes.
\end{exampleblock}
\end{frame}

\begin{frame}{O que é Biologia Sintética?}
\definicao{Biologia Sintética}{
Disciplina que combina biologia, engenharia e ciência da computação para projetar e construir novas funções e sistemas biológicos que não existem na natureza.
}

\vspace{0.3cm}

\begin{columns}
\column{0.5\textwidth}
\textbf{Princípios de Engenharia:}
\begin{itemize}
    \item Modularidade
    \item Padronização
    \item Abstração
    \item Design racional
    \item Otimização
\end{itemize}

\column{0.5\textwidth}
\textbf{Aplicações:}
\begin{itemize}
    \item Biocombustíveis
    \item Fármacos
    \item Biosensores
    \item Biomateriais
    \item Terapia celular
    \item Bioremediação
\end{itemize}
\end{columns}
\end{frame}

\begin{frame}{Ciclo Design-Build-Test-Learn}
\begin{center}
\begin{tikzpicture}[node distance=3cm, auto,
    box/.style={rectangle, draw, fill=azulclaro!20, text width=2.2cm, text centered, rounded corners, minimum height=1.2cm}]

    \node[box, fill=roxodna!20] (design) {DESIGN\\Modelar\\Projetar};
    \node[box, fill=verdeacido!20, right of=design] (build) {BUILD\\Construir\\Sintetizar};
    \node[box, fill=laranjacel!20, below of=build] (test) {TEST\\Testar\\Medir};
    \node[box, fill=azulclaro!30, left of=test] (learn) {LEARN\\Analisar\\Aprender};

    \draw[->, thick, azulescuro, line width=1.5pt] (design) -- (build);
    \draw[->, thick, azulescuro, line width=1.5pt] (build) -- (test);
    \draw[->, thick, azulescuro, line width=1.5pt] (test) -- (learn);
    \draw[->, thick, azulescuro, line width=1.5pt] (learn) -- (design);

\end{tikzpicture}
\end{center}

\importante{
Processo iterativo: cada ciclo refina o design com base nos dados experimentais
}
\end{frame}

\begin{frame}{Ética e Biossegurança}
\begin{columns}
\column{0.5\textwidth}
\begin{alertblock}{Considerações Éticas}
\begin{itemize}
    \item Organismos geneticamente modificados (OGMs)
    \item Impacto ambiental
    \item Uso dual (militar/civil)
    \item Bioterrorismo
    \item Propriedade intelectual
\end{itemize}
\end{alertblock}

\column{0.5\textwidth}
\begin{block}{Níveis de Biossegurança}
\begin{itemize}
    \item \textbf{BSL-1:} risco mínimo
    \item \textbf{BSL-2:} risco moderado
    \item \textbf{BSL-3:} alto risco
    \item \textbf{BSL-4:} risco extremo
\end{itemize}
\end{block}

\begin{exampleblock}{Regulação}
Protocolos de contenção, monitoramento e aprovação ética
\end{exampleblock}
\end{columns}
\end{frame}

% ============================================================================
% SEÇÃO 2: CIRCUITOS GENÉTICOS SINTÉTICOS
% ============================================================================
\section{Circuitos Genéticos Sintéticos}

\begin{frame}{Circuitos Genéticos como Dispositivos Computacionais}
\begin{block}{Analogia com Circuitos Eletrônicos}
Circuitos genéticos processam sinais bioquímicos assim como circuitos eletrônicos processam sinais elétricos
\end{block}

\vspace{0.3cm}

\begin{center}
\begin{tabular}{lcc}
\toprule
\textbf{Conceito} & \textbf{Eletrônica} & \textbf{Biologia} \\
\midrule
Sinal & Voltagem & Concentração \\
Componente & Transistor & Gene/Proteína \\
Estado ON & Alto voltagem & Alta expressão \\
Estado OFF & Baixo voltagem & Baixa expressão \\
Lógica & Portas lógicas & Regulação gênica \\
\bottomrule
\end{tabular}
\end{center}

\importante{
Princípios de design modular: componentes intercambiáveis e reutilizáveis
}
\end{frame}

\begin{frame}{BioBricks: Partes Biológicas Padronizadas}
\begin{columns}
\column{0.5\textwidth}
\begin{block}{Componentes Básicos}
\begin{enumerate}
    \item \textbf{Promotor:} controla transcrição
    \item \textbf{RBS:} Ribosome Binding Site
    \item \textbf{CDS:} Coding Sequence (gene)
    \item \textbf{Terminador:} fim da transcrição
\end{enumerate}
\end{block}

\textbf{Montagem Modular:}
\begin{itemize}
    \item BioBrick Assembly
    \item Gibson Assembly
    \item Golden Gate
\end{itemize}

\column{0.5\textwidth}
\begin{center}
\begin{tikzpicture}[scale=0.8]
    % Promotor
    \draw[->, thick, verdeacido, line width=2pt] (0,0) -- (1,0);
    \node[below, font=\tiny] at (0.5,-0.1) {Promotor};

    % RBS
    \draw[thick, roxodna, fill=roxodna!30] (1.2,0) circle (0.15);
    \node[below, font=\tiny] at (1.2,-0.3) {RBS};

    % CDS
    \draw[->, thick, azulescuro, line width=2pt] (1.5,0) -- (3.5,0);
    \node[below, font=\tiny] at (2.5,-0.1) {Gene (CDS)};

    % Terminador
    \draw[thick, laranjacel, line width=3pt] (3.7,-0.2) -- (3.7,0.2);
    \node[below, font=\tiny] at (3.7,-0.3) {Term};

    % Expressão
    \draw[->, thick, red, dashed] (2.5,0.5) -- (2.5,1) node[above, font=\small] {Proteína};
\end{tikzpicture}
\end{center}

\vspace{0.5cm}

\begin{exampleblock}{Registry of Standard Biological Parts}
\url{http://parts.igem.org}
\end{exampleblock}
\end{columns}
\end{frame}

\begin{frame}{Toggle Switch: Interruptor Genético}
\begin{columns}
\column{0.5\textwidth}
\textbf{Gardner et al. (2000)}

\begin{block}{Princípio}
Dois repressores que se inibem mutuamente
\begin{itemize}
    \item Gene A reprime gene B
    \item Gene B reprime gene A
    \item Sistema bistável
\end{itemize}
\end{block}

\textbf{Estados Estáveis:}
\begin{itemize}
    \item Estado 1: A alto, B baixo
    \item Estado 2: A baixo, B alto
\end{itemize}

\column{0.5\textwidth}
\begin{center}
\begin{tikzpicture}[scale=0.9]
    % Gene A
    \node[draw, rectangle, fill=verdeacido!30, minimum width=2cm, minimum height=0.8cm] (geneA) at (0,2) {Gene A};

    % Gene B
    \node[draw, rectangle, fill=laranjacel!30, minimum width=2cm, minimum height=0.8cm] (geneB) at (0,0) {Gene B};

    % Proteína A
    \node[draw, circle, fill=verdeacido!50, minimum size=0.6cm] (protA) at (3,2) {A};

    % Proteína B
    \node[draw, circle, fill=laranjacel!50, minimum size=0.6cm] (protB) at (3,0) {B};

    % Produção
    \draw[->, thick, verdeacido] (geneA) -- (protA);
    \draw[->, thick, laranjacel] (geneB) -- (protB);

    % Repressão cruzada
    \draw[-|, thick, red, bend left=30] (protA) to (geneB);
    \draw[-|, thick, red, bend right=30] (protB) to (geneA);

    \node[right, font=\tiny] at (4,2) {Repressor A};
    \node[right, font=\tiny] at (4,0) {Repressor B};
\end{tikzpicture}
\end{center}
\end{columns}

\importante{
Primeiro circuito genético sintético com memória bistável
}
\end{frame}

\begin{frame}{Modelo Matemático do Toggle Switch}
\begin{block}{Sistema de EDOs}
\begin{align*}
\derivada{u}{t} &= \frac{\alpha_1}{1 + v^{\beta}} - u\\
\derivada{v}{t} &= \frac{\alpha_2}{1 + u^{\gamma}} - v
\end{align*}

onde:
\begin{itemize}
    \item $u, v$ = concentrações dos repressores
    \item $\alpha_1, \alpha_2$ = taxas de produção máximas
    \item $\beta, \gamma$ = coeficientes de Hill (cooperatividade)
\end{itemize}
\end{block}

\textbf{Condição para Bistabilidade:}
\destaque{
\begin{equation*}
\alpha_1 \alpha_2 > (\beta + 1)(\gamma + 1)
\end{equation*}
}
\end{frame}

\begin{frame}{Análise de Bistabilidade}
\begin{columns}
\column{0.5\textwidth}
\textbf{Pontos de Equilíbrio:}
\begin{equation*}
u = \frac{\alpha_1}{1 + v^{\beta}}, \quad v = \frac{\alpha_2}{1 + u^{\gamma}}
\end{equation*}

\vspace{0.3cm}

\textbf{Análise de Estabilidade:}
\begin{itemize}
    \item Jacobiano
    \item Autovalores
    \item Diagrama de bifurcação
\end{itemize}

\vspace{0.3cm}

\begin{exampleblock}{Resultado}
Sistema pode ter 1 ou 3 pontos de equilíbrio dependendo dos parâmetros
\end{exampleblock}

\column{0.5\textwidth}
\begin{center}
\begin{tikzpicture}[scale=0.8]
    \draw[->] (0,0) -- (4,0) node[right] {$u$};
    \draw[->] (0,0) -- (0,4) node[above] {$v$};

    % Nullclines
    \draw[thick, verdeacido, domain=0:3.5, samples=100] plot (\x, {3/(1+\x^2)});
    \draw[thick, laranjacel, domain=0:3.5, samples=100] plot ({3/(1+\x^2)}, \x);

    % Pontos de equilíbrio
    \fill[blue] (0.5,2.8) circle (2pt) node[above right, font=\tiny] {Estável};
    \fill[red] (1.5,1.5) circle (2pt) node[above right, font=\tiny] {Instável};
    \fill[blue] (2.8,0.5) circle (2pt) node[above right, font=\tiny] {Estável};

    % Trajetórias
    \draw[->, thick, gray, dashed] (0.3,1.5) -- (0.4,2.5);
    \draw[->, thick, gray, dashed] (2.5,1.5) -- (2.6,0.7);

    \node[below, font=\small] at (2,-0.3) {Repressor A};
    \node[left, font=\small, rotate=90] at (-0.5,2) {Repressor B};
\end{tikzpicture}
\end{center}

\textbf{Nullclines:} $\derivada{u}{t} = 0$ e $\derivada{v}{t} = 0$
\end{columns}
\end{frame}

\begin{frame}{Repressilator: Oscilador Genético}
\textbf{Elowitz \& Leibler (2000)}

\begin{columns}
\column{0.5\textwidth}
\begin{block}{Princípio}
Três genes que se reprimem em cascata:
\begin{itemize}
    \item Gene \gene{lacI} reprime \gene{tetR}
    \item Gene \gene{tetR} reprime \gene{cI}
    \item Gene \gene{cI} reprime \gene{lacI}
\end{itemize}
\end{block}

\textbf{Resultado:}
\begin{itemize}
    \item Oscilações periódicas
    \item Período $\sim$ 2-3 horas
    \item Visualização com GFP
\end{itemize}

\column{0.5\textwidth}
\begin{center}
\begin{tikzpicture}[scale=1.0]
    % Três genes em círculo
    \node[draw, circle, fill=verdeacido!30, minimum size=1.2cm] (lacI) at (90:2) {\gene{lacI}};
    \node[draw, circle, fill=laranjacel!30, minimum size=1.2cm] (tetR) at (210:2) {\gene{tetR}};
    \node[draw, circle, fill=roxodna!30, minimum size=1.2cm] (cI) at (330:2) {\gene{cI}};

    % Repressão circular
    \draw[-|, thick, red, line width=1.5pt] (lacI) to[bend right=20] (tetR);
    \draw[-|, thick, red, line width=1.5pt] (tetR) to[bend right=20] (cI);
    \draw[-|, thick, red, line width=1.5pt] (cI) to[bend right=20] (lacI);

    % Centro
    \node[font=\small] at (0,0) {Ciclo de\\repressão};
\end{tikzpicture}
\end{center}

\vspace{0.3cm}

\importante{
Primeiro oscilador genético sintético
}
\end{columns}
\end{frame}

\begin{frame}{Modelo Matemático do Repressilator}
\begin{block}{Sistema de EDOs}
\begin{align*}
\derivada{m_i}{t} &= -m_i + \frac{\alpha}{1 + p_j^n} + \alpha_0\\
\derivada{p_i}{t} &= -\beta(p_i - m_i)
\end{align*}

onde $i = 1, 2, 3$ (lacI, tetR, cI) e $j = i-1 \pmod{3}$

\begin{itemize}
    \item $m_i$ = concentração de mRNA do gene $i$
    \item $p_i$ = concentração de proteína do gene $i$
    \item $\alpha$ = taxa de transcrição máxima
    \item $n$ = coeficiente de Hill
    \item $\beta$ = razão de degradação proteína/mRNA
\end{itemize}
\end{block}

\textbf{Condições para Oscilação:}
\begin{itemize}
    \item Repressão forte ($n$ grande)
    \item Tempo de resposta adequado ($\beta$ apropriado)
\end{itemize}
\end{frame}

\begin{frame}{Feed-Forward Loops (FFLs)}
\begin{columns}
\column{0.5\textwidth}
\begin{block}{Motivo de Rede}
Padrão comum em redes regulatórias:
\begin{itemize}
    \item Gene X regula Y
    \item Gene X regula Z
    \item Gene Y regula Z
\end{itemize}
\end{block}

\textbf{Tipos:}
\begin{itemize}
    \item Coherent FFL (C-FFL)
    \item Incoherent FFL (I-FFL)
\end{itemize}

\textbf{Funções:}
\begin{itemize}
    \item Filtro de ruído
    \item Detector de pulso
    \item Acelerador de resposta
\end{itemize}

\column{0.5\textwidth}
\begin{center}
\begin{tikzpicture}[scale=1.0]
    % Type 1 Coherent FFL
    \node[draw, circle, fill=verdeacido!30] (X) at (0,2) {X};
    \node[draw, circle, fill=laranjacel!30] (Y) at (2,2) {Y};
    \node[draw, circle, fill=roxodna!30] (Z) at (1,0) {Z};

    % Ativação
    \draw[->, thick, azulescuro, line width=1.5pt] (X) -- (Y);
    \draw[->, thick, azulescuro, line width=1.5pt] (X) to[bend right=20] (Z);
    \draw[->, thick, azulescuro, line width=1.5pt] (Y) to[bend left=20] (Z);

    \node[above, font=\small] at (1,2.5) {C1-FFL};
    \node[below, font=\tiny] at (1,-0.5) {AND gate};
\end{tikzpicture}

\vspace{0.5cm}

\begin{tikzpicture}[scale=1.0]
    % Type 1 Incoherent FFL
    \node[draw, circle, fill=verdeacido!30] (X) at (0,2) {X};
    \node[draw, circle, fill=laranjacel!30] (Y) at (2,2) {Y};
    \node[draw, circle, fill=roxodna!30] (Z) at (1,0) {Z};

    % Ativação/Repressão
    \draw[->, thick, azulescuro, line width=1.5pt] (X) -- (Y);
    \draw[->, thick, azulescuro, line width=1.5pt] (X) to[bend right=20] (Z);
    \draw[-|, thick, red, line width=1.5pt] (Y) to[bend left=20] (Z);

    \node[above, font=\small] at (1,2.5) {I1-FFL};
    \node[below, font=\tiny] at (1,-0.5) {Pulse detector};
\end{tikzpicture}
\end{center}
\end{columns}
\end{frame}

\begin{frame}{Portas Lógicas Genéticas}
\begin{center}
\begin{tikzpicture}[scale=0.9]
    % AND gate
    \begin{scope}
        \node[draw, rectangle, fill=azulclaro!30, minimum width=1.5cm, minimum height=1cm] (and) at (0,0) {AND};
        \draw[<-] (and.west) ++(-0.5,0.2) -- (and.west |- 0,0.2) node[left, font=\tiny] {A};
        \draw[<-] (and.west) ++(-0.5,-0.2) -- (and.west |- 0,-0.2) node[left, font=\tiny] {B};
        \draw[->] (and.east) -- ++(0.5,0) node[right, font=\tiny] {Out};
        \node[below, font=\small] at (0,-0.8) {A e B ativos};
    \end{scope}

    % OR gate
    \begin{scope}[xshift=4cm]
        \node[draw, rectangle, fill=verdeacido!30, minimum width=1.5cm, minimum height=1cm] (or) at (0,0) {OR};
        \draw[<-] (or.west) ++(-0.5,0.2) -- (or.west |- 0,0.2) node[left, font=\tiny] {A};
        \draw[<-] (or.west) ++(-0.5,-0.2) -- (or.west |- 0,-0.2) node[left, font=\tiny] {B};
        \draw[->] (or.east) -- ++(0.5,0) node[right, font=\tiny] {Out};
        \node[below, font=\small] at (0,-0.8) {A ou B ativo};
    \end{scope}

    % NOT gate
    \begin{scope}[xshift=8cm]
        \node[draw, rectangle, fill=laranjacel!30, minimum width=1.5cm, minimum height=1cm] (not) at (0,0) {NOT};
        \draw[<-] (not.west) ++(-0.5,0) -- (not.west) node[left, font=\tiny] {A};
        \draw[->] (not.east) -- ++(0.5,0) node[right, font=\tiny] {Out};
        \node[below, font=\small] at (0,-0.8) {Inversor};
    \end{scope}

    % NOR gate
    \begin{scope}[xshift=11.5cm]
        \node[draw, rectangle, fill=roxodna!30, minimum width=1.5cm, minimum height=1cm] (nor) at (0,0) {NOR};
        \draw[<-] (nor.west) ++(-0.5,0.2) -- (nor.west |- 0,0.2) node[left, font=\tiny] {A};
        \draw[<-] (nor.west) ++(-0.5,-0.2) -- (nor.west |- 0,-0.2) node[left, font=\tiny] {B};
        \draw[->] (nor.east) -- ++(0.5,0) node[right, font=\tiny] {Out};
        \node[below, font=\small] at (0,-0.8) {NOT(A OR B)};
    \end{scope}
\end{tikzpicture}
\end{center}

\vspace{0.5cm}

\begin{exampleblock}{Implementação Biológica}
\begin{itemize}
    \item \textbf{AND:} dois promotores necessários
    \item \textbf{OR:} promotores alternativos
    \item \textbf{NOT:} repressor
    \item \textbf{NOR:} dois repressores
\end{itemize}
\end{exampleblock}
\end{frame}

\begin{frame}{Biosensores e Repórteres}
\begin{columns}
\column{0.5\textwidth}
\begin{block}{Biosensor}
Sistema que detecta sinais químicos/físicos e produz resposta mensurável
\end{block}

\textbf{Componentes:}
\begin{itemize}
    \item Domínio sensor (receptor)
    \item Domínio de sinal
    \item Gene repórter
\end{itemize}

\vspace{0.3cm}

\textbf{Repórteres Comuns:}
\begin{itemize}
    \item GFP (fluorescência verde)
    \item RFP (fluorescência vermelha)
    \item Luciferase (bioluminescência)
    \item LacZ ($\beta$-galactosidase)
\end{itemize}

\column{0.5\textwidth}
\begin{center}
\begin{tikzpicture}[scale=0.85]
    % Molécula alvo
    \node[draw, circle, fill=red!30, minimum size=0.6cm] (target) at (0,3) {X};
    \node[above, font=\tiny] at (0,3.4) {Analito};

    % Sensor
    \node[draw, rectangle, fill=verdeacido!30, minimum width=2.5cm, minimum height=0.8cm] (sensor) at (0,2) {Promotor\\Sensor};

    % Arrow
    \draw[->, thick, red] (target) -- (sensor);

    % Gene repórter
    \draw[->, thick, azulescuro, line width=2pt] (0,1.2) -- (2,1.2);
    \node[below, font=\tiny] at (1,1.2) {GFP};

    % Terminador
    \draw[thick, laranjacel, line width=3pt] (2.2,1) -- (2.2,1.4);

    % Fluorescência
    \draw[->, thick, verdeacido, dashed] (1,0.8) -- (1,0.2);
    \node[draw, star, fill=verdeacido!50, star points=5, star point ratio=2, minimum size=0.8cm] at (1,-0.3) {};
    \node[below, font=\small] at (1,-0.8) {Fluorescência};
\end{tikzpicture}
\end{center}

\vspace{0.3cm}

\exemplo{Arsênico}{
Biosensor para detecção de arsênico em água usando promotor \textit{ars} + GFP
}
\end{columns}
\end{frame}

\begin{frame}[fragile]{Simulação: Toggle Switch em Python}
\begin{lstlisting}[language=Python, caption=Modelo do toggle switch, basicstyle=\ttfamily\tiny]
import numpy as np
import matplotlib.pyplot as plt
from scipy.integrate import odeint

def toggle_switch(y, t, alpha1, alpha2, beta, gamma):
    """Modelo do toggle switch"""
    u, v = y
    dudt = alpha1 / (1 + v**beta) - u
    dvdt = alpha2 / (1 + u**gamma) - v
    return [dudt, dvdt]

# Parametros para sistema bistavel
alpha1 = 156.25
alpha2 = 15.6
beta = 2.5
gamma = 1.0

# Tempo
t = np.linspace(0, 50, 1000)

# Duas condicoes iniciais diferentes
y0_1 = [0.1, 5.0]  # Estado 1
y0_2 = [5.0, 0.1]  # Estado 2

# Resolver
sol1 = odeint(toggle_switch, y0_1, t, args=(alpha1, alpha2, beta, gamma))
sol2 = odeint(toggle_switch, y0_2, t, args=(alpha1, alpha2, beta, gamma))

# Plotar
plt.figure(figsize=(10, 6))
plt.plot(t, sol1[:, 0], 'b-', label='u (Estado 1)', linewidth=2)
plt.plot(t, sol1[:, 1], 'r-', label='v (Estado 1)', linewidth=2)
plt.plot(t, sol2[:, 0], 'b--', label='u (Estado 2)', linewidth=2)
plt.plot(t, sol2[:, 1], 'r--', label='v (Estado 2)', linewidth=2)
plt.xlabel('Tempo')
plt.ylabel('Concentracao')
plt.legend()
plt.title('Toggle Switch: Bistabilidade')
\end{lstlisting}
\end{frame}

\begin{frame}[fragile]{Simulação: Repressilator em Python}
\begin{lstlisting}[language=Python, caption=Modelo do repressilator, basicstyle=\ttfamily\tiny]
def repressilator(y, t, alpha, alpha0, beta, n):
    """Modelo do repressilator (3 genes)"""
    m1, p1, m2, p2, m3, p3 = y

    # mRNA equations
    dm1dt = -m1 + alpha / (1 + p3**n) + alpha0
    dm2dt = -m2 + alpha / (1 + p1**n) + alpha0
    dm3dt = -m3 + alpha / (1 + p2**n) + alpha0

    # Protein equations
    dp1dt = -beta * (p1 - m1)
    dp2dt = -beta * (p2 - m2)
    dp3dt = -beta * (p3 - m3)

    return [dm1dt, dp1dt, dm2dt, dp2dt, dm3dt, dp3dt]

# Parametros
alpha = 216.0
alpha0 = 0.0
beta = 0.2
n = 2.0

# Condicoes iniciais
y0 = [0.5, 1.0, 1.5, 0.5, 2.0, 1.0]
t = np.linspace(0, 500, 5000)

# Resolver
sol = odeint(repressilator, y0, t, args=(alpha, alpha0, beta, n))

# Plotar
plt.figure(figsize=(12, 6))
plt.subplot(1, 2, 1)
plt.plot(t, sol[:, 1], 'g-', label='Proteina 1', linewidth=2)
plt.plot(t, sol[:, 3], 'orange', label='Proteina 2', linewidth=2)
plt.plot(t, sol[:, 5], 'purple', label='Proteina 3', linewidth=2)
plt.xlabel('Tempo')
plt.ylabel('Concentracao')
plt.legend()
plt.title('Repressilator: Series Temporais')

plt.subplot(1, 2, 2)
plt.plot(sol[:, 1], sol[:, 3], 'b-', linewidth=1)
plt.xlabel('Proteina 1')
plt.ylabel('Proteina 2')
plt.title('Retrato de Fase')
\end{lstlisting}
\end{frame}

% ============================================================================
% SEÇÃO 3: ENGENHARIA METABÓLICA
% ============================================================================
\section{Engenharia Metabólica}

\begin{frame}{Engenharia de Vias Metabólicas}
\definicao{Engenharia Metabólica}{
Modificação direcionada de vias metabólicas usando tecnologia de DNA recombinante para melhorar a produção de compostos de interesse.
}

\vspace{0.3cm}

\begin{columns}
\column{0.5\textwidth}
\textbf{Objetivos:}
\begin{itemize}
    \item Aumentar rendimento
    \item Produzir novos compostos
    \item Remover subprodutos
    \item Melhorar eficiência
    \item Reduzir custos
\end{itemize}

\column{0.5\textwidth}
\textbf{Alvos Típicos:}
\begin{itemize}
    \item Biocombustíveis (etanol, biodiesel)
    \item Fármacos (artemisinina, taxol)
    \item Aminoácidos
    \item Vitaminas
    \item Polímeros biodegradáveis
\end{itemize}
\end{columns}

\vspace{0.3cm}

\importante{
Combina biologia molecular, bioquímica e modelagem computacional
}
\end{frame}

\begin{frame}{Estratégias de Engenharia Metabólica}
\begin{block}{1. Overexpression (Superexpressão)}
Aumentar expressão de enzimas limitantes da via
\begin{itemize}
    \item Promotores fortes
    \item Múltiplas cópias do gene
    \item Otimização de códons
\end{itemize}
\end{block}

\begin{block}{2. Gene Knockout (Deleção)}
Eliminar vias competidoras ou indesejáveis
\begin{itemize}
    \item CRISPR-Cas9
    \item Recombinação homóloga
    \item Direcionar fluxo para via de interesse
\end{itemize}
\end{block}

\begin{block}{3. Heterologous Expression (Expressão Heteróloga)}
Introduzir genes de outros organismos
\begin{itemize}
    \item Criar novas vias
    \item Produzir compostos não-nativos
\end{itemize}
\end{block}
\end{frame}

\begin{frame}{Estratégia Push-Pull}
\begin{center}
\begin{tikzpicture}[scale=0.9, node distance=2cm]
    % Substrato
    \node[draw, circle, fill=azulclaro!30, minimum size=1cm] (S) at (0,0) {S};

    % Intermediário
    \node[draw, circle, fill=verdeacido!30, minimum size=1cm] (I) at (3,0) {I};

    % Produto
    \node[draw, circle, fill=laranjacel!30, minimum size=1cm] (P) at (6,0) {P};

    % Via competidora
    \node[draw, circle, fill=gray!30, minimum size=0.8cm] (X) at (3,-2) {X};

    % Setas principais
    \draw[->, thick, verdeacido, line width=2pt] (S) -- (I) node[midway, above, font=\small] {PUSH};
    \draw[->, thick, laranjacel, line width=2pt] (I) -- (P) node[midway, above, font=\small] {PULL};

    % Via competidora
    \draw[->, thick, gray] (I) -- (X);
    \draw[red, line width=3pt] (2.8,-0.8) -- (3.2,-1.2);
    \draw[red, line width=3pt] (3.2,-0.8) -- (2.8,-1.2);

    % Labels
    \node[above, font=\small] at (1.5,0.5) {Superexpressão};
    \node[above, font=\small] at (4.5,0.5) {Superexpressão};
    \node[below, font=\small] at (3,-2.5) {Knockout};

    \node[below, font=\tiny] at (0,-0.7) {Substrato};
    \node[below, font=\tiny] at (3,-0.7) {Intermediário};
    \node[below, font=\tiny] at (6,-0.7) {Produto alvo};
\end{tikzpicture}
\end{center}

\vspace{0.3cm}

\begin{exampleblock}{Estratégia}
\begin{itemize}
    \item \textbf{Push:} aumentar fluxo entrando na via
    \item \textbf{Pull:} aumentar conversão para produto final
    \item \textbf{Block:} eliminar vias competidoras
\end{itemize}
\end{exampleblock}
\end{frame}

\begin{frame}{Balanceamento de Cofatores}
\begin{block}{Problema}
Muitas reações requerem cofatores (NADH, NADPH, ATP)
\begin{itemize}
    \item Disponibilidade limitada
    \item Necessário balanço entre vias
    \item Afeta rendimento máximo teórico
\end{itemize}
\end{block}

\vspace{0.3cm}

\begin{center}
\begin{tikzpicture}[scale=0.9]
    % Ciclo NADH/NAD+
    \node[draw, circle, fill=verdeacido!30] (nadh) at (0,0) {NADH};
    \node[draw, circle, fill=laranjacel!30] (nad) at (3,0) {NAD$^+$};

    % Reação catabólica
    \draw[->, thick, azulescuro, bend left=45] (nad) to node[above, font=\small] {Catabolismo} (nadh);

    % Reação anabólica
    \draw[->, thick, roxodna, bend left=45] (nadh) to node[below, font=\small] {Anabolismo} (nad);

    % Labels
    \node[below, font=\small] at (1.5,-1.2) {Balanço redox};
\end{tikzpicture}
\end{center}

\importante{
Engenharia de cofatores: modificar especificidade de enzimas (NADH $\leftrightarrow$ NADPH)
}
\end{frame}

\begin{frame}{Exemplo: Produção de Artemisinina}
\begin{columns}
\column{0.5\textwidth}
\begin{block}{Contexto}
\begin{itemize}
    \item Antimalárico potente
    \item Extraído de \textit{Artemisia annua}
    \item Rendimento baixo ($\sim$1\%)
    \item Alto custo
\end{itemize}
\end{block}

\textbf{Solução: Biologia Sintética}
\begin{itemize}
    \item Via heteróloga em \textit{E. coli}
    \item Via estendida em levedura
    \item Expressão de genes de planta
    \item Otimização metabólica
\end{itemize}

\textbf{Resultado (Keasling lab):}
\begin{itemize}
    \item Produção \textgreater 25 g/L
    \item Redução de custo 90\%
\end{itemize}

\column{0.5\textwidth}
\begin{center}
\begin{tikzpicture}[scale=0.7, node distance=1.3cm]
    % Via do mevalonato
    \node[draw, rectangle, fill=azulclaro!30, text width=2cm, align=center] (acetyl) at (0,4) {Acetil-CoA};

    \node[draw, rectangle, fill=azulclaro!30, text width=2cm, align=center] (mev) at (0,2.5) {Mevalonato};

    \node[draw, rectangle, fill=verdeacido!30, text width=2cm, align=center] (ipp) at (0,1) {IPP/DMAPP};

    \node[draw, rectangle, fill=laranjacel!30, text width=2cm, align=center] (fpp) at (0,-0.5) {FPP};

    \node[draw, rectangle, fill=roxodna!30, text width=2cm, align=center] (amor) at (0,-2) {Amorfa-\\dieno};

    \node[draw, rectangle, fill=red!30, text width=2cm, align=center] (art) at (0,-3.5) {Ácido\\Artemisínico};

    % Setas
    \draw[->, thick] (acetyl) -- (mev) node[midway, right, font=\tiny] {ERG};
    \draw[->, thick] (mev) -- (ipp) node[midway, right, font=\tiny] {MVA};
    \draw[->, thick] (ipp) -- (fpp) node[midway, right, font=\tiny] {ERG20};
    \draw[->, thick, verdeacido, line width=1.5pt] (fpp) -- (amor) node[midway, right, font=\tiny] {ADS};
    \draw[->, thick, verdeacido, line width=1.5pt] (amor) -- (art) node[midway, right, font=\tiny] {CYP71AV1};

    % Labels
    \node[left, font=\tiny] at (-1.5,2.5) {Nativo};
    \node[left, font=\tiny, text=verdeacido] at (-1.5,-2.5) {Heterólogo};
\end{tikzpicture}
\end{center}
\end{columns}
\end{frame}

\begin{frame}{Exemplo: Produção de 1,3-Propanodiol}
\begin{columns}
\column{0.5\textwidth}
\begin{block}{Aplicação}
Monômero para polímeros (PTT)
\begin{itemize}
    \item Carpetes, plásticos
    \item Produção tradicional: petroquímica
    \item Alternativa: fermentação
\end{itemize}
\end{block}

\textbf{Via Sintética:}
\begin{itemize}
    \item Glicerol $\rightarrow$ 3-HPA
    \item 3-HPA $\rightarrow$ 1,3-PDO
    \item Genes de \textit{Klebsiella}
    \item Expressão em \textit{E. coli}
\end{itemize}

\textbf{DuPont/Genencor:}
\begin{itemize}
    \item Produção industrial
    \item Rendimento \textgreater 135 g/L
    \item Processo econômico
\end{itemize}

\column{0.5\textwidth}
\begin{center}
\begin{tikzpicture}[scale=0.8, node distance=1.5cm]
    % Glicerol
    \node[draw, rectangle, fill=azulclaro!30, minimum width=2cm] (gly) at (0,3) {Glicerol};

    % 3-HPA
    \node[draw, rectangle, fill=laranjacel!30, minimum width=2cm] (hpa) at (0,1.5) {3-HPA};

    % 1,3-PDO
    \node[draw, rectangle, fill=verdeacido!30, minimum width=2cm] (pdo) at (0,0) {1,3-PDO};

    % Via competidora
    \node[draw, rectangle, fill=gray!30, minimum width=1.5cm] (dha) at (3,1.5) {DHA};

    % Setas
    \draw[->, thick, verdeacido, line width=2pt] (gly) -- (hpa) node[midway, left, font=\small] {dhaB};
    \draw[->, thick, verdeacido, line width=2pt] (hpa) -- (pdo) node[midway, left, font=\small] {dhaT};

    % Via competidora
    \draw[->, thick, gray] (hpa) -- (dha);
    \draw[red, line width=3pt] (2.2,1.4) -- (2.6,1.6);
    \draw[red, line width=3pt] (2.6,1.4) -- (2.2,1.6);

    % Cofatores
    \node[right, font=\tiny] at (0.5,2.3) {NADH};
    \node[right, font=\tiny] at (0.5,0.8) {NADH};
\end{tikzpicture}
\end{center}

\vspace{0.5cm}

\importante{
Exemplo de sucesso comercial da engenharia metabólica
}
\end{columns}
\end{frame}

\begin{frame}{Análise de Balanço de Fluxo (FBA)}
\begin{block}{Flux Balance Analysis}
Método computacional para predizer fluxos metabólicos
\begin{itemize}
    \item Baseado em estequiometria
    \item Estado estacionário
    \item Otimização linear
\end{itemize}
\end{block}

\textbf{Formulação:}
\destaque{
\begin{align*}
\text{maximizar} \quad & c^T v\\
\text{sujeito a} \quad & Sv = 0\\
& v_{min} \leq v \leq v_{max}
\end{align*}
}

onde:
\begin{itemize}
    \item $S$ = matriz estequiométrica
    \item $v$ = vetor de fluxos
    \item $c$ = coeficientes da função objetivo (ex: biomassa)
\end{itemize}
\end{frame}

\begin{frame}{Identificação de Alvos Genéticos}
\begin{columns}
\column{0.5\textwidth}
\begin{block}{Métodos}
\begin{itemize}
    \item \textbf{FBA:} predizer fluxos
    \item \textbf{FVA:} análise de variabilidade
    \item \textbf{MOMA:} minimização de ajuste metabólico
    \item \textbf{OptKnock:} identificar knockouts ótimos
\end{itemize}
\end{block}

\textbf{OptKnock:}
\begin{itemize}
    \item Problema bi-nível
    \item Acoplar crescimento e produção
    \item Identificar deleções sinérgicas
\end{itemize}

\column{0.5\textwidth}
\begin{center}
\begin{tikzpicture}[scale=0.8]
    % Gráfico biomassa vs produto
    \draw[->] (0,0) -- (4.5,0) node[right, font=\small] {Biomassa};
    \draw[->] (0,0) -- (0,3.5) node[above, font=\small] {Produto};

    % Linha de compensação wild-type
    \draw[thick, gray, dashed] (0,0) -- (4,0.5) node[right, font=\tiny] {Wild-type};

    % Região factível após knockout
    \fill[verdeacido!30, opacity=0.5] (0,0) -- (1.5,2.5) -- (3,2) -- (4,1) -- (4,0) -- cycle;

    % Ponto ótimo
    \fill[red] (1.5,2.5) circle (2pt);
    \node[above, font=\small] at (1.5,2.7) {Ótimo};

    % Legenda
    \node[below, font=\tiny, text width=3cm, align=center] at (2,-0.7) {OptKnock identifica knockouts para produção acoplada};
\end{tikzpicture}
\end{center}
\end{columns}
\end{frame}

\begin{frame}{Engenharia de Cofatores}
\begin{block}{Estratégias}
\begin{itemize}
    \item Modificar especificidade NADH/NADPH
    \item Superexpressar transidrogenases
    \item Redirecionar fluxo de cofatores
    \item Usar vias alternativas
\end{itemize}
\end{block}

\vspace{0.3cm}

\begin{columns}
\column{0.5\textwidth}
\textbf{Exemplo: NADPH}
\begin{itemize}
    \item Necessário para biossíntese
    \item Via das pentoses (PPP)
    \item Enzima málica
    \item Transidrogenase PntAB
\end{itemize}

\column{0.5\textwidth}
\begin{center}
\begin{tikzpicture}[scale=0.7]
    % NADP+ e NADPH
    \node[draw, circle, fill=laranjacel!30] (nadp) at (0,2) {NADP$^+$};
    \node[draw, circle, fill=verdeacido!30] (nadph) at (0,0) {NADPH};

    % NAD+ e NADH
    \node[draw, circle, fill=laranjacel!30] (nad) at (3,2) {NAD$^+$};
    \node[draw, circle, fill=verdeacido!30] (nadh) at (3,0) {NADH};

    % Transidorgenase
    \draw[->, thick, azulescuro, bend left=20] (nadh) to node[above, font=\tiny] {PntAB} (nadph);
    \draw[->, thick, azulescuro, bend left=20] (nadp) to (nad);

    % PPP
    \draw[->, thick, roxodna] (nadp) to[out=180, in=180, looseness=2] node[left, font=\tiny] {PPP} (nadph);
\end{tikzpicture}
\end{center}
\end{columns}

\importante{
Balanço de cofatores é crítico para rendimento máximo
}
\end{frame}

\begin{frame}{Regulação Dinâmica de Vias}
\begin{block}{Controle Temporal}
Expressão gênica responsiva a sinais metabólicos
\begin{itemize}
    \item Promotores induzíveis
    \item Biosensores de metabólitos
    \item Controle de feedback
    \item Evitar toxicidade de intermediários
\end{itemize}
\end{block}

\vspace{0.3cm}

\begin{center}
\begin{tikzpicture}[scale=0.9]
    % Fase de crescimento
    \draw[thick, verdeacido, line width=2pt] (0,0) -- (2,2) node[midway, above, sloped, font=\small] {Crescimento};

    % Fase de produção
    \draw[thick, laranjacel, line width=2pt] (2,2) -- (5,2) node[midway, above, font=\small] {Produção};

    % Eixos
    \draw[->] (0,0) -- (5.5,0) node[right, font=\small] {Tempo};
    \draw[->] (0,0) -- (0,2.5) node[above, font=\small] {Biomassa/Produto};

    % Divisão
    \draw[dashed] (2,0) -- (2,2.5);
    \node[below, font=\small] at (2,-0.2) {Switch};

    % Labels
    \node[below, font=\tiny] at (1,-0.5) {Via de biomassa ativa};
    \node[below, font=\tiny] at (3.5,-0.5) {Via de produto ativa};
\end{tikzpicture}
\end{center}

\exemplo{Sistema Bifásico}{
Fase 1: crescimento celular (promotor constitutivo)\\
Fase 2: produção (promotor induzível por IPTG ou temperatura)
}
\end{frame}

\begin{frame}[fragile]{Engenharia Metabólica com COBRApy}
\begin{lstlisting}[language=Python, caption=Análise FBA com COBRApy, basicstyle=\ttfamily\tiny]
import cobra
from cobra import Model, Reaction, Metabolite

# Carregar modelo de E. coli
model = cobra.io.load_json_model("e_coli_core.json")

# Realizar FBA
solution = model.optimize()
print(f"Taxa de crescimento: {solution.objective_value:.3f} h^-1")

# Analisar fluxos
print("\nFluxos principais:")
for rxn_id in ['BIOMASS_Ecoli_core_w_GAM', 'EX_glc__D_e', 'EX_o2_e']:
    rxn = model.reactions.get_by_id(rxn_id)
    print(f"{rxn.name}: {solution.fluxes[rxn_id]:.3f}")

# Simular knockout
model_ko = model.copy()
model_ko.reactions.PGI.knock_out()  # Knockout phosphoglucose isomerase

solution_ko = model_ko.optimize()
print(f"\nCrescimento pos-knockout: {solution_ko.objective_value:.3f} h^-1")

# Analise de variabilidade de fluxo (FVA)
from cobra.flux_analysis import flux_variability_analysis

fva_result = flux_variability_analysis(model, reaction_list=['PGI', 'PFK'])
print("\nFVA:")
print(fva_result)

# Identificar genes essenciais
from cobra.flux_analysis import single_gene_deletion

deletion_results = single_gene_deletion(model)
essential_genes = deletion_results[deletion_results['growth'] < 0.01]
print(f"\nNumero de genes essenciais: {len(essential_genes)}")
\end{lstlisting}
\end{frame}

\begin{frame}[fragile]{FSEOF: Flux Scanning com Objetivo Forçado}
\begin{lstlisting}[language=Python, caption=FSEOF para engenharia metabólica, basicstyle=\ttfamily\tiny]
def fseof_analysis(model, target_reaction, biomass_fraction=0.1):
    """
    Flux Scanning based on Enforced Objective Flux
    Identifica reacoes para superexpressao/delecao
    """
    import pandas as pd

    # Forcar producao minima de biomassa
    biomass_rxn = model.reactions.BIOMASS_Ecoli_core_w_GAM
    max_biomass = model.optimize().objective_value
    biomass_rxn.bounds = (biomass_fraction * max_biomass, 1000)

    # Maximizar produto alvo
    model.objective = target_reaction

    results = []

    # Escanear cada reacao
    for rxn in model.reactions:
        if rxn.id == target_reaction.id:
            continue

        # Estado original
        original_bounds = rxn.bounds

        # Testar knockout
        rxn.bounds = (0, 0)
        sol_ko = model.optimize()
        flux_ko = sol_ko.objective_value if sol_ko.status == 'optimal' else 0

        # Testar superexpressao (aumentar limite superior)
        rxn.bounds = (original_bounds[0], original_bounds[1] * 10)
        sol_oe = model.optimize()
        flux_oe = sol_oe.objective_value if sol_oe.status == 'optimal' else 0

        # Restaurar
        rxn.bounds = original_bounds

        results.append({
            'reaction': rxn.id,
            'knockout_effect': flux_ko,
            'overexpression_effect': flux_oe
        })

    return pd.DataFrame(results)

# Executar FSEOF
target = model.reactions.EX_succ_e  # Succinato como produto alvo
fseof_results = fseof_analysis(model, target)

# Identificar alvos promissores
print("Melhores alvos para superexpressao:")
print(fseof_results.nlargest(5, 'overexpression_effect'))
\end{lstlisting}
\end{frame}

% ============================================================================
% SEÇÃO 4: FERRAMENTAS DE DESIGN E MODELAGEM
% ============================================================================
\section{Ferramentas de Design e Modelagem}

\begin{frame}{Ferramentas CAD para Biologia Sintética}
\begin{block}{Computer-Aided Design (CAD)}
Software para projetar circuitos e vias biológicas
\end{block}

\vspace{0.3cm}

\begin{columns}
\column{0.5\textwidth}
\textbf{Ferramentas Principais:}
\begin{itemize}
    \item \textbf{Cello:} design automático de circuitos
    \item \textbf{j5:} montagem de DNA
    \item \textbf{Eugene:} linguagem de design
    \item \textbf{SBOLDesigner:} editor visual
    \item \textbf{iBioSim:} modelagem e simulação
\end{itemize}

\column{0.5\textwidth}
\textbf{Funcionalidades:}
\begin{itemize}
    \item Design de construtos
    \item Predição de comportamento
    \item Otimização de parâmetros
    \item Seleção de partes
    \item Planejamento de montagem
\end{itemize}
\end{columns}

\vspace{0.3cm}

\begin{exampleblock}{Cello}
Entrada: especificação lógica (Verilog)\\
Saída: sequência de DNA otimizada
\end{exampleblock}
\end{frame}

\begin{frame}{Caracterização e Padronização de Partes}
\begin{block}{Caracterização}
Medição quantitativa de propriedades de partes biológicas
\begin{itemize}
    \item Força de promotor (PoPS - Polymerase per Second)
    \item Eficiência de RBS (taxa de tradução)
    \item Taxa de degradação
    \item Resposta a indutores
\end{itemize}
\end{block}

\vspace{0.3cm}

\begin{columns}
\column{0.5\textwidth}
\textbf{Registry of Standard Biological Parts}
\begin{itemize}
    \item \textgreater 20,000 partes
    \item Caracterização padronizada
    \item Compartilhamento comunitário
    \item iGEM competition
\end{itemize}

\column{0.5\textwidth}
\begin{center}
\begin{tikzpicture}[scale=0.8]
    % Eixos
    \draw[->] (0,0) -- (4,0) node[right, font=\small] {[Indutor]};
    \draw[->] (0,0) -- (0,3) node[above, font=\small] {Expressão};

    % Curva dose-resposta
    \draw[thick, verdeacido, domain=0:3.5, samples=100] plot (\x, {2.5/(1+exp(-2*(\x-1.5)))});

    % Parâmetros
    \draw[dashed] (0,1.25) -- (1.5,1.25) -- (1.5,0);
    \node[left, font=\tiny] at (0,1.25) {$K_{1/2}$};
    \node[below, font=\tiny] at (1.5,0) {EC$_{50}$};
    \draw[dashed] (0,2.5) -- (3.8,2.5);
    \node[left, font=\tiny] at (0,2.5) {Max};
\end{tikzpicture}
\end{center}

\textbf{Curva dose-resposta}
\end{columns}
\end{frame}

\begin{frame}{SBOL: Synthetic Biology Open Language}
\definicao{SBOL}{
Formato de dados padronizado para representar designs de biologia sintética
}

\vspace{0.3cm}

\begin{columns}
\column{0.5\textwidth}
\textbf{Objetivos:}
\begin{itemize}
    \item Interoperabilidade entre ferramentas
    \item Compartilhamento de designs
    \item Documentação estruturada
    \item Reprodutibilidade
\end{itemize}

\textbf{Componentes:}
\begin{itemize}
    \item ComponentDefinition
    \item ModuleDefinition
    \item Sequence
    \item Interaction
\end{itemize}

\column{0.5\textwidth}
\begin{center}
\begin{tikzpicture}[scale=0.7]
    % Diagrama hierárquico
    \node[draw, rectangle, fill=azulclaro!30, text width=2.5cm, align=center] (design) at (0,3) {Design};

    \node[draw, rectangle, fill=verdeacido!30, text width=1.8cm, align=center] (module) at (-1.5,1.5) {Módulo};
    \node[draw, rectangle, fill=laranjacel!30, text width=1.8cm, align=center] (comp) at (1.5,1.5) {Componente};

    \node[draw, rectangle, fill=roxodna!30, text width=1.5cm, align=center] (seq) at (0,0) {Sequência};

    \draw[->, thick] (design) -- (module);
    \draw[->, thick] (design) -- (comp);
    \draw[->, thick] (module) -- (seq);
    \draw[->, thick] (comp) -- (seq);
\end{tikzpicture}
\end{center}

\vspace{0.5cm}

\textbf{Versão atual: SBOL 3.0}
\end{columns}

\importante{
Padrão aceito pela comunidade para intercâmbio de dados
}
\end{frame}

\begin{frame}{Frameworks de Modelagem}
\begin{columns}
\column{0.5\textwidth}
\begin{block}{Modelos Determinísticos}
\begin{itemize}
    \item \textbf{ODEs:} sistemas contínuos
    \item \textbf{Steady-state:} FBA
    \item Ferramentas: COPASI, Tellurium
\end{itemize}
\end{block}

\begin{block}{Modelos Estocásticos}
\begin{itemize}
    \item \textbf{Gillespie:} simulação exata
    \item \textbf{Tau-leaping:} aproximado
    \item Para baixo número de moléculas
\end{itemize}
\end{block}

\column{0.5\textwidth}
\begin{block}{Modelos Espaciais}
\begin{itemize}
    \item \textbf{PDEs:} difusão e reação
    \item \textbf{Agent-based:} células individuais
    \item Virtual Cell
\end{itemize}
\end{block}

\begin{block}{Modelos Multi-escala}
\begin{itemize}
    \item Integrar múltiplos níveis
    \item Molecular $\rightarrow$ Celular $\rightarrow$ População
    \item CompuCell3D
\end{itemize}
\end{block}
\end{columns}

\vspace{0.3cm}

\begin{exampleblock}{Escolha do Modelo}
Depende da escala, número de moléculas e questão biológica
\end{exampleblock}
\end{frame}

\begin{frame}{Algoritmos de Otimização para Design}
\begin{block}{OptKnock}
Identifica knockouts genéticos ótimos
\begin{itemize}
    \item Problema bi-nível
    \item Otimização mista inteira (MILP)
    \item Acopla crescimento e produção
\end{itemize}
\end{block}

\begin{block}{OptGene}
Extensão usando algoritmos genéticos
\begin{itemize}
    \item Busca heurística
    \item Múltiplos objetivos
    \item Mais flexível que OptKnock
\end{itemize}
\end{block}

\begin{block}{OptForce}
Identifica intervenções mínimas
\begin{itemize}
    \item Superexpressão e knockout
    \item Garante produção mínima
\end{itemize}
\end{block}
\end{frame}

\begin{frame}{Machine Learning para Predição de Partes}
\begin{columns}
\column{0.5\textwidth}
\begin{block}{Aplicações de ML}
\begin{itemize}
    \item Predição de força de promotor
    \item Otimização de RBS
    \item Predição de estrutura de RNA
    \item Design de proteínas
    \item Predição de interações
\end{itemize}
\end{block}

\textbf{Métodos:}
\begin{itemize}
    \item Regressão
    \item Redes neurais
    \item Random forests
    \item Deep learning
\end{itemize}

\column{0.5\textwidth}
\begin{center}
\begin{tikzpicture}[scale=0.8]
    % Rede neural simples
    % Input layer
    \foreach \i in {1,2,3} {
        \node[draw, circle, fill=azulclaro!30, minimum size=0.5cm] (i\i) at (0,\i) {};
    }
    \node[above, font=\small] at (0,3.5) {Input};

    % Hidden layer
    \foreach \i in {1,2,3,4} {
        \node[draw, circle, fill=verdeacido!30, minimum size=0.5cm] (h\i) at (2,\i-0.5) {};
    }
    \node[above, font=\small] at (2,3.5) {Hidden};

    % Output layer
    \node[draw, circle, fill=laranjacel!30, minimum size=0.5cm] (o1) at (4,2) {};
    \node[above, font=\small] at (4,2.7) {Output};

    % Connections
    \foreach \i in {1,2,3} {
        \foreach \j in {1,2,3,4} {
            \draw[->, thin, gray] (i\i) -- (h\j);
        }
    }
    \foreach \i in {1,2,3,4} {
        \draw[->, thin, gray] (h\i) -- (o1);
    }

    \node[below, font=\tiny] at (0,0.5) {Sequência};
    \node[below, font=\tiny] at (4,1.5) {Expressão};
\end{tikzpicture}
\end{center}

\vspace{0.3cm}

\importante{
ML acelera ciclo de design-build-test-learn
}
\end{columns}
\end{frame}

\begin{frame}[fragile]{Simulação Estocástica: Algoritmo de Gillespie}
\begin{lstlisting}[language=Python, caption=Gillespie SSA para circuitos genéticos, basicstyle=\ttfamily\tiny]
import numpy as np
import matplotlib.pyplot as plt

def gillespie_ssa(reactions, initial_state, tmax):
    """
    Gillespie Stochastic Simulation Algorithm
    reactions: lista de (propensity_func, stoichiometry_change)
    """
    state = np.array(initial_state, dtype=float)
    time = 0
    times = [time]
    states = [state.copy()]

    while time < tmax:
        # Calcular propensidades
        propensities = np.array([prop(state) for prop, _ in reactions])
        a0 = propensities.sum()

        if a0 == 0:
            break

        # Tempo ate proxima reacao
        tau = np.random.exponential(1/a0)

        # Escolher reacao
        r = np.random.random() * a0
        cumsum = 0
        for i, a in enumerate(propensities):
            cumsum += a
            if cumsum >= r:
                _, change = reactions[i]
                state += change
                break

        time += tau
        times.append(time)
        states.append(state.copy())

    return np.array(times), np.array(states)

# Exemplo: producao e degradacao de proteina
k_prod = 10.0  # taxa de producao
k_deg = 0.5    # taxa de degradacao

reactions = [
    (lambda s: k_prod, np.array([1])),        # Producao
    (lambda s: k_deg * s[0], np.array([-1]))  # Degradacao
]

times, states = gillespie_ssa(reactions, [0], tmax=20)
plt.plot(times, states[:, 0])
plt.xlabel('Tempo')
plt.ylabel('Numero de moleculas')
\end{lstlisting}
\end{frame}

\begin{frame}[fragile]{Otimização de Vias com Algoritmos Genéticos}
\begin{lstlisting}[language=Python, caption=AG para otimização de expressão gênica, basicstyle=\ttfamily\tiny]
import numpy as np
from scipy.integrate import odeint

def pathway_model(y, t, gene_expression_levels):
    """Modelo de via metabolica simples"""
    S, I, P = y
    k1, k2 = gene_expression_levels  # Niveis de expressao

    dSdt = -k1 * S
    dIdt = k1 * S - k2 * I
    dPdt = k2 * I

    return [dSdt, dIdt, dPdt]

def fitness_function(gene_levels):
    """Objetivo: maximizar produto final"""
    y0 = [10, 0, 0]  # [S, I, P]
    t = np.linspace(0, 10, 100)
    sol = odeint(pathway_model, y0, t, args=(gene_levels,))
    return sol[-1, 2]  # Concentracao final de P

def genetic_algorithm(pop_size=50, generations=100):
    """Algoritmo genetico simples"""
    # Populacao inicial
    population = np.random.uniform(0, 10, (pop_size, 2))

    for gen in range(generations):
        # Avaliar fitness
        fitness = np.array([fitness_function(ind) for ind in population])

        # Selecao
        sorted_idx = np.argsort(fitness)[::-1]
        population = population[sorted_idx]

        # Elitismo: manter melhores 20%
        n_elite = pop_size // 5
        new_pop = population[:n_elite].copy()

        # Crossover e mutacao
        while len(new_pop) < pop_size:
            parents = population[np.random.choice(n_elite, 2, replace=False)]
            child = (parents[0] + parents[1]) / 2  # Crossover
            child += np.random.normal(0, 0.5, 2)   # Mutacao
            child = np.clip(child, 0, 10)
            new_pop = np.vstack([new_pop, child])

        population = new_pop[:pop_size]

    return population[0]

# Executar
optimal_levels = genetic_algorithm()
print(f"Niveis otimos de expressao: {optimal_levels}")
\end{lstlisting}
\end{frame}

% ============================================================================
% SEÇÃO 5: BIOLOGIA SINTÉTICA DE CÉLULAS INTEIRAS
% ============================================================================
\section{Biologia Sintética de Células Inteiras}

\begin{frame}{Genomas Mínimos}
\begin{columns}
\column{0.5\textwidth}
\begin{block}{Conceito}
Conjunto mínimo de genes necessários para vida
\begin{itemize}
    \item Redução sistemática do genoma
    \item Identificar genes essenciais
    \item Criar chassis simplificado
\end{itemize}
\end{block}

\textbf{Vantagens:}
\begin{itemize}
    \item Menor carga metabólica
    \item Mais previsível
    \item Reduz interações indesejadas
    \item Plataforma para engenharia
\end{itemize}

\column{0.5\textwidth}
\exemplo{JCVI-syn3.0}{
\begin{itemize}
    \item \textit{Mycoplasma mycoides}
    \item 531 kb (vs 1079 kb original)
    \item 473 genes
    \item Primeiro genoma mínimo sintético
    \item Venter Institute, 2016
\end{itemize}
}

\vspace{0.3cm}

\begin{alertblock}{Desafios}
\begin{itemize}
    \item Crescimento lento
    \item Genes de função desconhecida
    \item Dependência de contexto
\end{itemize}
\end{alertblock}
\end{columns}
\end{frame}

\begin{frame}{Xenobiologia e Código Genético Expandido}
\begin{columns}
\column{0.5\textwidth}
\definicao{Xenobiologia}{
Criação de sistemas biológicos com química alternativa à base da vida natural
}

\textbf{Abordagens:}
\begin{itemize}
    \item Bases não-naturais (X, Y)
    \item Aminoácidos sintéticos
    \item XNA (xeno nucleic acids)
    \item Ribossomos ortogonais
\end{itemize}

\textbf{Aplicações:}
\begin{itemize}
    \item Contenção biológica
    \item Propriedades novas
    \item Bioortogonalidade
\end{itemize}

\column{0.5\textwidth}
\begin{center}
\begin{tikzpicture}[scale=0.9]
    % DNA natural
    \node[draw, rectangle, fill=azulclaro!30, text width=2.5cm, align=center] (nat) at (0,2.5) {DNA Natural\\A, T, G, C};

    % DNA expandido
    \node[draw, rectangle, fill=verdeacido!30, text width=2.5cm, align=center] (exp) at (0,0.8) {DNA Expandido\\A, T, G, C, X, Y};

    % Seta
    \draw[->, thick, laranjacel, line width=2pt] (nat) -- (exp);
    \node[right, font=\small] at (0.5,1.6) {Expansão};

    % Aplicações
    \node[draw, rectangle, fill=roxodna!30, text width=2.5cm, align=center] (app) at (0,-0.8) {Proteínas com\\aminoácidos\\não-naturais};

    \draw[->, thick] (exp) -- (app);
\end{tikzpicture}
\end{center}

\vspace{0.3cm}

\importante{
Romesberg lab: organismos com 6 letras genéticas (2017)
}
\end{columns}
\end{frame}

\begin{frame}{Sistemas Cell-Free}
\begin{block}{Síntese Cell-Free}
Produção de proteínas e metabolismo sem células vivas
\begin{itemize}
    \item Extratos celulares
    \item Componentes purificados (PURE system)
    \item Adicionar DNA/RNA template
\end{itemize}
\end{block}

\vspace{0.3cm}

\begin{columns}
\column{0.5\textwidth}
\textbf{Vantagens:}
\begin{itemize}
    \item Sem viabilidade celular
    \item Controle preciso
    \item Sem membranas
    \item Prototipagem rápida
    \item Proteínas tóxicas
\end{itemize}

\column{0.5\textwidth}
\textbf{Aplicações:}
\begin{itemize}
    \item Produção de proteínas
    \item Biossensores
    \item Terapêutica
    \item Educação
    \item Biologia sintética
\end{itemize}
\end{columns}

\vspace{0.3cm}

\begin{exampleblock}{Produtos Comerciais}
\begin{itemize}
    \item PURExpress (NEB)
    \item myTXTL (Arbor Biosciences)
\end{itemize}
\end{exampleblock}
\end{frame}

\begin{frame}{Protocélulas e Células Artificiais}
\begin{columns}
\column{0.5\textwidth}
\begin{block}{Protocélulas}
Estruturas sintéticas que mimetizam propriedades celulares
\begin{itemize}
    \item Vesículas lipídicas
    \item Contêm maquinaria molecular
    \item Auto-replicação primitiva
    \item Origem da vida
\end{itemize}
\end{block}

\textbf{Componentes:}
\begin{itemize}
    \item Membrana (lipossomas)
    \item Genoma (DNA/RNA)
    \item Metabolismo (enzimas)
    \item Informação hereditária
\end{itemize}

\column{0.5\textwidth}
\begin{center}
\begin{tikzpicture}[scale=1.0]
    % Membrana
    \draw[thick, azulescuro, fill=azulclaro!10] (0,0) circle (1.5);

    % DNA
    \draw[thick, roxodna, decorate, decoration={snake, amplitude=0.3mm, segment length=2mm}] (-0.8,-0.3) -- (0.8,0.3);
    \node[below, font=\tiny, text=roxodna] at (0,-0.5) {DNA};

    % Ribossomos
    \fill[verdeacido] (-0.5,0.5) circle (1.5pt);
    \fill[verdeacido] (0.3,0.7) circle (1.5pt);
    \fill[verdeacido] (0.6,-0.4) circle (1.5pt);
    \node[above, font=\tiny, text=verdeacido] at (0,0.9) {Ribossomos};

    % Enzimas
    \draw[thick, laranjacel] (-0.6,-0.7) -- (-0.4,-0.9);
    \draw[thick, laranjacel] (0.5,0.2) circle (2pt);
    \node[below, font=\tiny, text=laranjacel] at (0.5,-1.0) {Enzimas};

    % Label
    \node[above, font=\small] at (0,1.8) {Protocélula};
\end{tikzpicture}
\end{center}

\vspace{0.3cm}

\importante{
Objetivo de longo prazo: criar vida artificial de novo
}
\end{columns}
\end{frame}

\begin{frame}{Sistemas Biológicos Ortogonais}
\definicao{Ortogonalidade}{
Sistemas que operam independentemente dos sistemas naturais da célula hospedeira
}

\vspace{0.3cm}

\begin{columns}
\column{0.5\textwidth}
\textbf{Exemplos:}
\begin{itemize}
    \item Ribossomos ortogonais
    \item tRNAs ortogonais
    \item Aminoacil-tRNA sintetases
    \item Códons não-naturais
    \item Fatores de transcrição ortogonais
\end{itemize}

\textbf{Vantagens:}
\begin{itemize}
    \item Sem interferência com host
    \item Controle independente
    \item Segurança biológica
\end{itemize}

\column{0.5\textwidth}
\begin{center}
\begin{tikzpicture}[scale=0.8]
    % Sistema natural
    \draw[thick, azulclaro, fill=azulclaro!10] (-1.5,1.5) rectangle (1.5,3.5);
    \node[above, font=\small] at (0,3.5) {Sistema Natural};

    \node[font=\tiny] at (0,2.5) {Ribossomo};
    \draw[->, thick] (-0.8,2.2) -- (-0.8,1.8) node[below, font=\tiny] {Proteína};

    % Sistema ortogonal
    \draw[thick, verdeacido, fill=verdeacido!10] (-1.5,-0.5) rectangle (1.5,1.5);
    \node[above, font=\small] at (0,1.5) {Sistema Ortogonal};

    \node[font=\tiny] at (0,0.5) {O-Ribossomo};
    \draw[->, thick] (0.8,0.2) -- (0.8,-0.2) node[below, font=\tiny] {O-Proteína};

    % Sem interferência
    \draw[red, thick, dashed] (-1.5,1.5) -- (1.5,1.5);
    \node[right, font=\tiny, text=red] at (1.6,1.5) {Isolamento};
\end{tikzpicture}
\end{center}

\vspace{0.3cm}

\exemplo{Aplicação}{
Incorporar aminoácidos não-naturais sem afetar proteoma endógeno
}
\end{columns}
\end{frame}

\begin{frame}{Aplicações: Células Terapêuticas}
\begin{columns}
\column{0.5\textwidth}
\begin{block}{Imunoterapia CAR-T}
Células T modificadas para atacar câncer
\begin{itemize}
    \item Receptor quimérico (CAR)
    \item Reconhece antígeno tumoral
    \item Resposta imune direcionada
    \item FDA aprovado (2017)
\end{itemize}
\end{block}

\textbf{Melhorias Sintéticas:}
\begin{itemize}
    \item Switches suicidas
    \item Controle temporal
    \item Multi-targeting
    \item Resistência à exaustão
\end{itemize}

\column{0.5\textwidth}
\begin{center}
\begin{tikzpicture}[scale=0.8]
    % Célula T
    \draw[thick, azulclaro, fill=azulclaro!20] (0,0) circle (1.2);
    \node[font=\small] at (0,0.5) {CAR-T};

    % CAR
    \draw[thick, verdeacido, line width=2pt] (-0.8,0.8) -- (-1.3,1.5);
    \node[above, font=\tiny, text=verdeacido] at (-1.3,1.6) {CAR};

    % Célula tumoral
    \draw[thick, red, fill=red!20] (3.5,0) circle (1.0);
    \node[font=\small] at (3.5,0) {Tumor};

    % Antígeno
    \draw[thick, red, fill=red!50] (2.5,0) circle (0.2);
    \node[below, font=\tiny] at (2.5,-0.4) {Antígeno};

    % Reconhecimento
    \draw[->, thick, verdeacido, line width=2pt] (-1.3,1.5) to[out=0, in=120] (2.5,0.2);

    % Morte celular
    \node[below, font=\small] at (3.5,-1.5) {Morte celular};
    \draw[->, thick, red, dashed] (3.5,-0.8) -- (3.5,-1.2);
\end{tikzpicture}
\end{center}

\vspace{0.3cm}

\importante{
Biologia sintética na medicina personalizada
}
\end{columns}
\end{frame}

\begin{frame}{Aplicações: Living Sensors}
\begin{block}{Biosensores Vivos}
Células modificadas para detectar e responder a sinais ambientais
\end{block}

\vspace{0.3cm}

\begin{columns}
\column{0.5\textwidth}
\textbf{Aplicações:}
\begin{itemize}
    \item Detecção de poluentes
    \item Diagnóstico de doenças
    \item Monitoramento intestinal
    \item Detecção de explosivos
    \item Controle de qualidade
\end{itemize}

\textbf{Componentes:}
\begin{itemize}
    \item Sensor (promotor responsivo)
    \item Processamento (circuito)
    \item Repórter (GFP, luz, etc)
    \item Memória (opcional)
\end{itemize}

\column{0.5\textwidth}
\begin{center}
\begin{tikzpicture}[scale=0.85]
    % Molécula alvo
    \node[draw, circle, fill=red!30, minimum size=0.6cm] (target) at (0,3) {Pb$^{2+}$};

    % Bactéria
    \draw[thick, azulclaro, fill=azulclaro!10, rounded corners] (-1.2,0.5) rectangle (1.2,2);
    \node[above, font=\small] at (0,2.2) {Bactéria sensora};

    % Circuito genético
    \draw[->, thick, verdeacido, line width=1.5pt] (-0.8,1.5) -- (0,1.5);
    \node[below, font=\tiny] at (-0.4,1.5) {Promotor};

    \draw[->, thick, roxodna, line width=1.5pt] (0.2,1.5) -- (0.8,1.5);
    \node[below, font=\tiny] at (0.5,1.5) {GFP};

    % Sinal de entrada
    \draw[->, thick, red] (target) -- (0,2);

    % Sinal de saída
    \draw[->, thick, verdeacido, dashed] (0,0.5) -- (0,0);
    \node[draw, star, fill=verdeacido!50, star points=5, minimum size=0.6cm] at (0,-0.5) {};
    \node[below, font=\small] at (0,-1) {Fluorescência};
\end{tikzpicture}
\end{center}
\end{columns}

\exemplo{Probiotic Sensor}{
\textit{E. coli} engineered para detectar inflamação intestinal e liberar terapêuticos
}
\end{frame}

% ============================================================================
% SEÇÃO 6: DESAFIOS E PERSPECTIVAS
% ============================================================================
\section{Desafios e Perspectivas}

\begin{frame}{Context Dependence e Genetic Burden}
\begin{columns}
\column{0.5\textwidth}
\begin{alertblock}{Context Dependence}
Comportamento de partes biológicas depende do contexto
\begin{itemize}
    \item Efeitos de posição
    \item Disponibilidade de recursos
    \item Carga de expressão
    \item Crosstalk regulatório
\end{itemize}
\end{alertblock}

\textbf{Problemas:}
\begin{itemize}
    \item Comportamento imprevisível
    \item Não-composicionalidade
    \item Necessidade de re-caracterização
\end{itemize}

\column{0.5\textwidth}
\begin{alertblock}{Genetic Burden}
Custo metabólico da expressão heteróloga
\begin{itemize}
    \item Ribossomos
    \item RNA polimerases
    \item Precursores (aminoácidos, NTPs)
    \item Energia (ATP, GTP)
\end{itemize}
\end{alertblock}

\textbf{Consequências:}
\begin{itemize}
    \item Crescimento lento
    \item Baixa produtividade
    \item Instabilidade
    \item Seleção contra circuito
\end{itemize}
\end{columns}

\vspace{0.3cm}

\importante{
Design deve considerar carga metabólica e recursos compartilhados
}
\end{frame}

\begin{frame}{Estabilidade Evolutiva}
\begin{block}{Problema}
Circuitos sintéticos podem ser perdidos ou modificados por evolução
\begin{itemize}
    \item Pressão seletiva contra carga metabólica
    \item Mutações inativadoras
    \item Perda de plasmídeos
    \item Rearranjos cromossomais
\end{itemize}
\end{block}

\vspace{0.3cm}

\begin{columns}
\column{0.5\textwidth}
\textbf{Estratégias de Estabilização:}
\begin{itemize}
    \item Integração cromossomal
    \item Acoplamento essencial
    \item Contenção genética
    \item Seleção positiva
    \item Toxina-antitoxina
\end{itemize}

\column{0.5\textwidth}
\begin{center}
\begin{tikzpicture}[scale=0.8]
    % Eixos
    \draw[->] (0,0) -- (4,0) node[right, font=\small] {Tempo};
    \draw[->] (0,0) -- (0,3) node[above, font=\small] {Função};

    % Circuito instável
    \draw[thick, red, domain=0:3.5] plot (\x, {2.5*exp(-0.5*\x)});
    \node[right, font=\tiny, text=red] at (3.5,0.5) {Instável};

    % Circuito estável
    \draw[thick, verdeacido, domain=0:3.5] plot (\x, {2});
    \node[right, font=\tiny, text=verdeacido] at (3.5,2) {Estável};
\end{tikzpicture}
\end{center}

\vspace{0.5cm}

\importante{
Trade-off entre função e fitness evolutivo
}
\end{columns}
\end{frame}

\begin{frame}{Escalonamento: Do Laboratório à Indústria}
\begin{block}{Desafios de Scale-up}
\begin{itemize}
    \item Condições de fermentação diferentes
    \item Gradientes de nutrientes e oxigênio
    \item Heterogeneidade populacional
    \item Controle de processo
    \item Custo de produção
    \item Regulação e aprovação
\end{itemize}
\end{block}

\vspace{0.3cm}

\begin{columns}
\column{0.5\textwidth}
\textbf{Etapas:}
\begin{enumerate}
    \item Shake flask (mL)
    \item Biorreator bancada (L)
    \item Piloto (10-100 L)
    \item Industrial (10,000+ L)
\end{enumerate}

\column{0.5\textwidth}
\textbf{Considerações:}
\begin{itemize}
    \item Mistura e aeração
    \item Controle pH e T
    \item Feeding strategies
    \item Downstream processing
    \item Purificação
    \item Custo de meios
\end{itemize}
\end{columns}

\vspace{0.3cm}

\exemplo{Sucesso}{
Produção de artemisinina e 1,3-PDO em escala industrial
}
\end{frame}

\begin{frame}{Questões Regulatórias e Éticas}
\begin{columns}
\column{0.5\textwidth}
\begin{alertblock}{Regulação}
\begin{itemize}
    \item Organismos Geneticamente Modificados (OGM)
    \item Protocolos de biossegurança
    \item Aprovação de produtos (FDA, EMA)
    \item Liberação ambiental
    \item Rastreabilidade
\end{itemize}
\end{alertblock}

\textbf{Níveis de Contenção:}
\begin{itemize}
    \item Física (laboratório)
    \item Biológica (auxotrofia)
    \item Genética (switches suicidas)
\end{itemize}

\column{0.5\textwidth}
\begin{alertblock}{Questões Éticas}
\begin{itemize}
    \item "Brincar de Deus"
    \item Impacto ecológico
    \item Bioterrorismo
    \item Acesso e equidade
    \item Propriedade intelectual
    \item Consentimento informado
\end{itemize}
\end{alertblock}

\begin{exampleblock}{Diretrizes}
\begin{itemize}
    \item Protocolo de Cartagena
    \item Convenção de Armas Biológicas
    \item iGEM Safety Committee
\end{itemize}
\end{exampleblock}
\end{columns}
\end{frame}

\begin{frame}{Biossegurança e Dual-Use}
\begin{block}{Dual-Use Research}
Pesquisa com potencial para uso benéfico e maléfico
\begin{itemize}
    \item Síntese de patógenos
    \item Engenharia de virulência
    \item Resistência a antibióticos
    \item Toxinas
\end{itemize}
\end{block}

\vspace{0.3cm}

\begin{columns}
\column{0.5\textwidth}
\textbf{Medidas de Segurança:}
\begin{itemize}
    \item Screening de síntese de DNA
    \item Controle de acesso
    \item Revisão ética
    \item Educação em biossegurança
    \item Códigos de conduta
\end{itemize}

\column{0.5\textwidth}
\textbf{Casos Históricos:}
\begin{itemize}
    \item Síntese de poliovírus (2002)
    \item Influenza 1918 (2005)
    \item H5N1 transmissível (2011)
\end{itemize}

\vspace{0.3cm}

\importante{
Balanço entre progresso científico e segurança pública
}
\end{columns}
\end{frame}

\begin{frame}{Direções Futuras}
\begin{columns}
\column{0.5\textwidth}
\begin{block}{AI-Driven Design}
\begin{itemize}
    \item Machine learning para predição
    \item Design generativo
    \item Otimização automática
    \item Análise de big data
    \item Modelagem multi-escala
\end{itemize}
\end{block}

\begin{block}{Foundries Automatizadas}
\begin{itemize}
    \item Robótica
    \item High-throughput
    \item Design-build-test automático
    \item Cloud labs
    \item Redução de custos e tempo
\end{itemize}
\end{block}

\column{0.5\textwidth}
\begin{block}{Novas Aplicações}
\begin{itemize}
    \item Medicina regenerativa
    \item Agricultura sustentável
    \item Materiais vivos
    \item Computação biológica
    \item Terraformação
    \item Exploração espacial
\end{itemize}
\end{block}

\begin{exampleblock}{Visão}
Biologia sintética como engenharia madura com design preditivo e confiável
\end{exampleblock}
\end{columns}

\vspace{0.3cm}

\importante{
Convergência de biologia, engenharia, computação e IA
}
\end{frame}

% ============================================================================
% SEÇÃO 7: EXERCÍCIOS
% ============================================================================
\section{Exercícios}

\begin{frame}{Exercícios - Parte 1}
\begin{block}{Questão 1: Análise de Bistabilidade}
Considere o toggle switch com parâmetros $\alpha_1 = 100$, $\alpha_2 = 20$, $\beta = 2$, $\gamma = 2$.

\begin{enumerate}
    \item Encontre os pontos de equilíbrio numericamente
    \item Calcule o Jacobiano em cada ponto
    \item Determine a estabilidade
    \item Verifique a condição de bistabilidade: $\alpha_1 \alpha_2 > (\beta + 1)(\gamma + 1)$
    \item Simule o sistema para diferentes condições iniciais
\end{enumerate}
\end{block}

\begin{block}{Questão 2: Design de Biosensor}
Projete um biosensor para detectar glicose:
\begin{enumerate}
    \item Identifique promotor responsivo a glicose
    \item Escolha gene repórter apropriado
    \item Desenhe o construto genético (BioBrick format)
    \item Estime a sensibilidade e faixa dinâmica
\end{enumerate}
\end{block}
\end{frame}

\begin{frame}{Exercícios - Parte 2}
\begin{block}{Questão 3: Otimização Metabólica}
Para produção de succinato em \textit{E. coli}:

\begin{enumerate}
    \item Use COBRApy para carregar modelo \textit{E. coli} core
    \item Maximize produção de succinato via FBA
    \item Identifique genes para knockout usando FSEOF
    \item Simule o efeito de knockouts
    \item Calcule rendimento teórico máximo (g succinato / g glicose)
\end{enumerate}
\end{block}

\begin{alertblock}{Dica}
Succinato: \texttt{EX\_succ\_e}\\
Glicose uptake: \texttt{EX\_glc\_\_D\_e}
\end{alertblock}
\end{frame}

\begin{frame}{Exercícios - Parte 3}
\begin{block}{Questão 4: Simulação Estocástica}
Implemente simulação de Gillespie para repressilator:

\begin{enumerate}
    \item Defina reações de transcrição e tradução
    \item Inclua degradação de mRNA e proteína
    \item Simule para 500 unidades de tempo
    \item Compare com solução determinística (ODEs)
    \item Analise variabilidade célula-a-célula
\end{enumerate}
\end{block}

\begin{block}{Questão 5: Design de Circuito Lógico}
Projete um circuito genético que implementa porta NAND:

\begin{enumerate}
    \item Desenhe diagrama do circuito
    \item Especifique componentes (promotores, genes, repressores)
    \item Construa tabela verdade
    \item Estime parâmetros necessários
    \item Discuta possíveis problemas de implementação
\end{enumerate}
\end{block}
\end{frame}

\begin{frame}{Projeto Integrador}
\begin{alertblock}{Projeto: Engenharia de Via para Produção de Bioplástico}
Objetivo: Produzir poli-3-hidroxibutirato (PHB) em \textit{E. coli}

\textbf{Tarefas:}
\begin{enumerate}
    \item \textbf{Design:}
    \begin{itemize}
        \item Identificar genes necessários (phaA, phaB, phaC)
        \item Projetar construtos com promotores apropriados
        \item Planejar estratégia de integração/plasmídeo
    \end{itemize}

    \item \textbf{Modelagem:}
    \begin{itemize}
        \item Construir modelo ODE da via
        \item Usar FBA para predizer rendimento
        \item Identificar gargalos metabólicos
    \end{itemize}

    \item \textbf{Otimização:}
    \begin{itemize}
        \item Propor knockouts para aumentar fluxo
        \item Balancear cofatores (NADPH)
        \item Simular estratégias de regulação dinâmica
    \end{itemize}

    \item \textbf{Análise:}
    \begin{itemize}
        \item Calcular rendimento teórico máximo
        \item Estimar custos de produção
        \item Discutir viabilidade comercial
    \end{itemize}
\end{enumerate}
\end{alertblock}
\end{frame}

% ============================================================================
% SEÇÃO 8: REFERÊNCIAS
% ============================================================================
\section{Referências}

\begin{frame}{Referências Principais}
\begin{thebibliography}{99}
\bibitem{gardner2000} Gardner, T.S., Cantor, C.R., \& Collins, J.J. (2000). Construction of a genetic toggle switch in \textit{Escherichia coli}. \textit{Nature}, 403, 339-342.

\bibitem{elowitz2000} Elowitz, M.B. \& Leibler, S. (2000). A synthetic oscillatory network of transcriptional regulators. \textit{Nature}, 403, 335-338.

\bibitem{keasling2010} Keasling, J.D. (2010). Manufacturing molecules through metabolic engineering. \textit{Science}, 330, 1355-1358.

\bibitem{nielsen2013} Nielsen, J. \& Keasling, J.D. (2016). Engineering cellular metabolism. \textit{Cell}, 164, 1185-1197.

\bibitem{cameron2014} Cameron, D.E., Bashor, C.J., \& Collins, J.J. (2014). A brief history of synthetic biology. \textit{Nature Reviews Microbiology}, 12, 381-390.
\end{thebibliography}
\end{frame}

\begin{frame}{Leitura Complementar}
\begin{block}{Livros}
\begin{itemize}
    \item Channon, K., Bromley, E.H., \& Woolfson, D.N. (2008). \textit{Synthetic Biology}. Oxford University Press.
    \item Smolke, C.D. (2018). \textit{The Metabolic Pathway Engineering Handbook}. CRC Press.
    \item Keasling, J.D. \& Venter, J.C. (2012). \textit{Synthetic Biology}. Academic Press.
\end{itemize}
\end{block}

\begin{block}{Reviews}
\begin{itemize}
    \item Khalil, A.S. \& Collins, J.J. (2010). Synthetic biology: applications come of age. \textit{Nature Reviews Genetics}, 11, 367-379.
    \item Church, G.M. et al. (2014). Realizing the potential of synthetic biology. \textit{Nature Reviews Molecular Cell Biology}, 15, 289-294.
    \item Lim, W.A. (2010). Designing customized cell signalling circuits. \textit{Nature Reviews Molecular Cell Biology}, 11, 393-403.
\end{itemize}
\end{block}
\end{frame}

\begin{frame}{Recursos e Ferramentas}
\begin{block}{Software e Ferramentas}
\begin{itemize}
    \item \textbf{COBRApy:} \url{https://opencobra.github.io/cobrapy/}
    \item \textbf{Tellurium:} \url{http://tellurium.analogmachine.org}
    \item \textbf{COPASI:} \url{http://copasi.org}
    \item \textbf{SBOLDesigner:} \url{https://github.com/SynBioDex/SBOLDesigner}
    \item \textbf{Cello:} \url{http://www.cellocad.org}
\end{itemize}
\end{block}

\begin{block}{Repositórios de Partes}
\begin{itemize}
    \item \textbf{iGEM Registry:} \url{http://parts.igem.org}
    \item \textbf{Addgene:} \url{https://www.addgene.org}
    \item \textbf{JBEI Registry:} \url{https://public-registry.jbei.org}
\end{itemize}
\end{block}

\begin{block}{Cursos Online}
\begin{itemize}
    \item MIT OpenCourseWare: Synthetic Biology
    \item Coursera: Bioinformatics Specialization
\end{itemize}
\end{block}
\end{frame}

% ============================================================================
% SLIDE FINAL
% ============================================================================
\begin{frame}
\begin{center}
{\Huge Obrigado!}

\vspace{1cm}

{\Large Capítulo 14: Design of Biological Systems}

\vspace{1cm}

\textbf{Tópicos Abordados:}\\
Circuitos genéticos sintéticos, engenharia metabólica,\\
ferramentas de design, células inteiras, desafios e perspectivas

\vspace{1cm}

\textit{Dúvidas? Perguntas?}
\end{center}
\end{frame}

\end{document}
